\section{Projective directional charts}

\subsection{Finite $z$ chart}

The finite chart represents the characteristic direction by
\[
r(z)=\binom{z}{1}.
\]
The fiber is parametrized by $(\tau,z)$ as established in the previous section.
For visualization only, the current application also uses
\[
\widehat z=\frac{2}{\pi}\arctan z,
\]
which compactifies finite $z$ into $(-1,1)$. This display coordinate is not a replacement for the projective chart.

\subsection{Reciprocal chart}

Near the direction $z=\infty$, introduce
\[
Z=\frac1z.
\]
Multiplying the affine fiber equation by $Z^2$ gives
\[
\widetilde P(Z)u+\widetilde Q(Z)v+\widetilde R(Z)=0,
\]
with
\begin{align*}
\widetilde P(Z)&=2b_{20}+(b_{11}-2a_{20})Z-a_{11}Z^2,\\
\widetilde Q(Z)&=b_{11}+(2b_{02}-a_{11})Z-2a_{02}Z^2,\\
\widetilde R(Z)&=b_{10}+(b_{01}-a_{10})Z-a_{01}Z^2.
\end{align*}
On the overlap $Z\neq0$,
\[
\widetilde P=Z^2P,\qquad
\widetilde Q=Z^2Q,\qquad
\widetilde R=Z^2R.
\]
The orthogonal foot $U_0$ is therefore unchanged, while the tangent vector scales as
\[
(-\widetilde Q,\widetilde P)=Z^2(-Q,P).
\]
If $\eta$ denotes the fiber coefficient in the reciprocal chart, equality of the physical midpoint requires
\[
U_0+\tau(-Q,P)
 =U_0+\eta(-\widetilde Q,\widetilde P),
\]
and hence
\[
\boxed{\eta=\frac{\tau}{Z^2}=\tau z^2},
\qquad
\boxed{\tau=\eta Z^2}.
\]
This transition follows from the general quadratic model and the chosen fiber parametrization; it is not imported as an assumption from the earlier wave-manifold application.

\paragraph{Full wave-state transition.}
The preceding derivation concerns the midpoint fiber coordinate.  For a full
wave state away from $\mathcal C$, the jump coordinate must transform as well.
Writing the reciprocal direction as $(1,Z)$, one has
\[
 (z,1)=Z^{-1}(1,Z),
\]
so the reciprocal jump coordinate is
\[
 \boxed{\Upsilon=Y/Z}.
\]
Thus the full reciprocal coordinates are $(\eta,\Upsilon,Z)$.  On
$\mathcal C$ both $Y$ and $\Upsilon$ vanish, which is why the distinction did
not affect the previously validated rarefaction calculations.
